\chapter*{ภาคผนวก ก: ค่าคงที่ทางฟิสิกส์มูลฐานสากล}
\addcontentsline{toc}{chapter}{ภาคผนวก ก: ค่าคงที่ทางฟิสิกส์มูลฐานสากล}

ตารางต่อไปนี้แสดงค่าคงที่ทางฟิสิกส์มูลฐานตามข้อแนะนำของคณะกรรมการข้อมูลสำหรับวิทยาศาสตร์และเทคโนโลยี (CODATA 2026 Recommended Values)

\vspace{0.4cm}
\begin{table}[htbp]
\centering
\small
\begin{tabularx}{\textwidth}{p{4.8cm} c l X}
\toprule
\textbf{ปริมาณทางฟิสิกส์} & \textbf{สัญลักษณ์} & \textbf{ค่าเชิงตัวเลข} & \textbf{หน่วยมาตรฐานเอสไอ} \\
\midrule
ความเร็วแสงในสุญญากาศ & $c$ & $299\,792\,458$ & $\text{m}\cdot\text{s}^{-1}$ (ค่าคงที่แม่นตรง) \\
ค่าคงที่พลังค์ & $h$ & $6.626\,070\,15 \times 10^{-34}$ & $\text{J}\cdot\text{s}$ (ค่าคงที่แม่นตรง) \\
ค่าคงที่พลังค์ลดทอน & $\hbar = h/2\pi$ & $1.054\,571\,817 \times 10^{-34}$ & $\text{J}\cdot\text{s}$ \\
ประจุธาตุ & $e$ & $1.602\,176\,634 \times 10^{-19}$ & $\text{C}$ (ค่าคงที่แม่นตรง) \\
ค่าคงที่โบลต์ซมันน์ & $k_B$ & $1.380\,649 \times 10^{-23}$ & $\text{J}\cdot\text{K}^{-1}$ (ค่าคงที่แม่นตรง) \\
เลขอาโวกาโดร & $N_A$ & $6.022\,140\,76 \times 10^{23}$ & $\text{mol}^{-1}$ (ค่าคงที่แม่นตรง) \\
ค่าคงที่โน้มถ่วงสากล & $G$ & $6.674\,30(15) \times 10^{-11}$ & $\text{m}^3\cdot\text{kg}^{-1}\cdot\text{s}^{-2}$ \\
สภาพยอมของสุญญากาศ & $\epsilon_0$ & $8.854\,187\,8128(13) \times 10^{-12}$ & $\text{F}\cdot\text{m}^{-1}$ \\
สภาพให้ซึมได้ของสุญญากาศ & $\mu_0$ & $1.256\,637\,062\,12(19) \times 10^{-6}$ & $\text{N}\cdot\text{A}^{-2}$ \\
มวลนิ่งของอิเล็กตรอน & $m_e$ & $9.109\,383\,7015(28) \times 10^{-31}$ & $\text{kg}$ \\
มวลนิ่งของโปรตอน & $m_p$ & $1.672\,621\,923\,69(51) \times 10^{-27}$ & $\text{kg}$ \\
มวลนิ่งของนิวตรอน & $m_n$ & $1.674\,927\,498\,04(95) \times 10^{-27}$ & $\text{kg}$ \\
รัศมีของบอร์ & $a_0$ & $5.291\,772\,109\,03(80) \times 10^{-11}$ & $\text{m}$ \\
\bottomrule
\end{tabularx}
\caption{ค่าคงที่ทางฟิสิกส์มูลฐานตามมาตรฐาน CODATA}
\label{tab:codata_constants}
\end{table}
